\section{Measuring the Effective History Span}
\label{sec:measure}

\paragraph{Setup.}
We consider multivariate forecasting with a history window $x_{t-L+1:t}$ and a horizon of $H=96$ steps, following the standard protocol \citep{dlinear, itransformer}. A lookback sweep varies $L \in \{96, 336, 720, 1440, 2880\}$. The confound we remove first is that longer windows reduce the number of available training windows, so any degradation could be attributed to data quantity rather than window content. Our protocol holds the number of training windows exactly constant at the count available under the longest lookback, subsampling uniformly for all shorter lookbacks (Appendix \ref{app:protocol}). We sweep four model families that differ in how history enters the computation, iTransformer \citep{itransformer} (attention across variates), PatchTST \citep{patchtst} (attention across patches), TimesNet \citep{timesnet} (2D convolutions over folded periods), and DLinear \citep{dlinear} (per-channel linear maps), on eight benchmarks (ETTh1, ETTh2, ETTm1, ETTm2, exchange rate, weather, electricity, traffic), three seeds each (the complete run ledger, including initially failed arms and their re-runs, is documented in Appendix \ref{app:ledger}).

\begin{table}[t]
\caption{Test MSE under the equal-data-budget protocol (iTransformer, three seeds per cell; full matrix in Appendix \ref{app:full}). Bold marks the best lookback within each row. The optimal lookback varies 30x across datasets.}
\label{tab:main}
\centering\scriptsize
\begin{tabular}{lcccccc}
\toprule
Dataset & $L{=}96$ & $L{=}336$ & $L{=}720$ & $L{=}1440$ & $L{=}2880$ & Best \\
\midrule
ETTh1 & 0.399 & 0.402 & \textbf{0.392} & 0.400 & 0.457 & 720 \\
ETTh2 & \textbf{0.299} & 0.302 & 0.305 & 0.325 & 0.445 & 96 \\
ETTm1 & 0.342 & \textbf{0.304} & 0.313 & 0.336 & 0.341 & 336 \\
ETTm2 & 0.185 & \textbf{0.174} & 0.179 & 0.187 & 0.191 & 336 \\
Electricity & 0.149 & 0.133 & 0.134 & 0.131 & \textbf{0.131} & 2880 \\
Exchange & \textbf{0.089} & 0.110 & 0.123 & 0.166 & 0.316 & 96 \\
Traffic & 0.406 & 0.365 & 0.358 & 0.348 & \textbf{0.341} & 2880 \\
Weather & 0.175 & \textbf{0.163} & 0.177 & 0.195 & 0.225 & 336 \\
\bottomrule
\end{tabular}
\end{table}


\paragraph{The divergence is real and large.}
Table \ref{tab:main} shows the controlled sweep for iTransformer. With the data budget held fixed, exchange rate degrades monotonically from 0.089 at $L{=}96$ to 0.316 at $L{=}2880$ (3.6x), while traffic improves monotonically from 0.406 to 0.341 ($-$16\%). Seed noise is far smaller than the lookback effect: the median per-seed standard deviation is 1.0\% of the cell mean across the full 480-run matrix, and the largest deviations occur in already-degraded long-window cells, where they change no ranking (the full distribution is in Appendix \ref{app:full}). Extending the sweep to $L{=}5760$ on electricity and traffic confirms that the benefit plateaus between $L{=}1440$ and $L{=}2880$ rather than being clipped by the scan range: errors at $L{=}5760$ never improve over $L{=}2880$ (Appendix \ref{app:full}). Table \ref{tab:bestl} summarizes the best lookback per model family: the data-side ordering---exchange shortest, electricity and traffic longest---is consistent across iTransformer, DLinear, PatchTST, and the zero-shot foundation model Chronos-Bolt \citep{chronos} (Section \ref{sec:zero}). TimesNet is the exception that exposes the model-side gate: its convolutional mixing cannot exploit long context, so its optimum stays at $L{=}96$ on electricity and traffic where every attention-based family goes long.

\begin{table}[t]
\caption{Best lookback per dataset and model family (equal data budget; all trained cells carry three seeds). $\uparrow$ marks monotone improvement through the longest span tested in that row; `--' marks settings not evaluated. $^{\dagger}$Wide-data zero-shot cells are evaluated with test-stride 2, Chronos-Bolt truncates contexts beyond 2048 steps, and electricity is suspected in-domain for the original Chronos corpus while traffic is confirmed never seen. The data-side ordering is model-consistent; the model side gates whether long context helps at all.}
\label{tab:bestl}
\centering\scriptsize
\begin{tabular}{lccccc}
\toprule
Dataset & iTransformer & DLinear & PatchTST & TimesNet & Chronos (zero-shot) \\
\midrule
ETTh1 & 720 & 1440 & 336 & 96 & $\uparrow$1440 \\
ETTh2 & 96 & 720 & 96 & 96 & 1440 \\
ETTm1 & 336 & 336 & 336 & 96 & $\uparrow$1440 \\
ETTm2 & 336 & 2880 & 336 & 96 & $\uparrow$1440 \\
Electricity & 2880 & 2880 & 1440 & 96 & 1440$^{\dagger}$ \\
Exchange & 96 & 96 & 96 & 96 & 96 \\
Traffic & 2880 & 2880 & $\uparrow$2880 & 96 & $\uparrow$2880$^{\dagger}$ \\
Weather & 336 & 2880 & 720 & 336 & $\uparrow$1440 \\
\bottomrule
\end{tabular}
\end{table}


\paragraph{It is not training dynamics.}
\label{sec:zero}
Chronos-Bolt, evaluated zero-shot with no training on our corpora, shows the same qualitative trend at the two poles of the ordering: on the short-span side it prefers $L{=}96$ on exchange rate; on the long-span side it improves monotonically on traffic up to the longest tested span (MSE 1.225 $\to$ 0.372, $-70\%$), and on electricity up to $L{=}1440$ ($-41\%$). On ETTh2 and weather the foundation model prefers longer spans ($L{=}1440$) than the trained families ($L{=}96$ and $L{=}336$); Appendix \ref{app:gift} attributes this gap to its strong per-instance normalization, which keeps the drift-harm channel dormant up to much higher drift than in the trained models. Bolt's context is capped at 2048 steps, so our $L{=}2880$ zero-shot runs read 2048 effective steps, and the wide-data zero-shot cells are evaluated with test-stride 2 (Appendix \ref{app:ledger}); while traffic is confirmed never seen in pretraining, electricity appears in the original Chronos corpus, so its zero-shot numbers carry a suspected-in-domain qualification and the long-span replication rests on traffic. The phenomenon is therefore a property of how the data's structure matches the inductive bias of sequence models in general, not of a specific training run. The ordering is also stable across horizons: at $H{=}192$, the short-, flat-, and long-span classes replicate exactly on exchange, ETTh1, and electricity (Appendix \ref{app:robust}). The two-axis classification further replicates on six additional public datasets (solar, PEMS04, PEMS08, ERCOT settlement prices, Melbourne pedestrian counts, and Mexico City bike demand; Appendix \ref{app:ext}): five are periodicity-dominated with long or saturating optimal spans, and the highest-drift newcomer (ERCOT) is the only one whose iTransformer optimum is short---its calendar statistic, the largest in the suite at 0.90, saturates at one week and does not pay back a longer window.

\paragraph{Definition.}
Let $e(L)$ be the test error of a fixed model class trained under the equal-data-budget protocol with lookback $L$. We define the \emph{effective history span} as
\begin{equation}
\mathrm{EHS} = \min\Big\{L : e(L) \le \min_{L'} e(L') \cdot (1 + \varepsilon)\Big\},
\label{eq:ehs}
\end{equation}
the smallest lookback whose error is within a fraction $\varepsilon$ of the best. This parsimonious form rewards shorter, cheaper windows once returns saturate. In the main text we report the argmin span $L^* = \arg\min_L e(L)$, which is the $\varepsilon \to 0$ limit and the quantity behind every ``best lookback'' figure; the parsimonious EHS($\varepsilon$) and its sensitivity to $\varepsilon$ are reported in Appendix \ref{app:robust}. The two coincide on drift-dominated datasets with a sharp optimum (exchange, ETTh2), and differ on benefit-dominated ones with a flat error curve near the optimum (electricity: $L^*{=}2880$ but EHS($0.02$)${=}336$).
